%% =======================================================================================
%% 
%% RELATÓRIO TÉCNICO
%% Criado a partir da classe abntex2, Modelo de Relatório Técnico do desenvolvido pelo 
%% Centro de Pesquisa em Arquitetura da Informação (CPAI) e Adaptado para a Coordenação de Curso Superior de Engenharia de Computação do IFFluminense Campus Bom Jesus do Itabapoana.
%% 
%% Arquivos de dados:   sty/configs.sty         % Configuração dos pacotes LaTeX
%%                      sty/metadados_rt.sty    % METADADOS DO RELATÓRIO - PREENCHER!
%%                      sty/dados_ccsec.sty      % Informações institucionais CCSEC / FUB
%%                      sty/timbrado.sty        % Formatação do documento (layout)
%%                      sty/vocabulario.sty     % Comandos LaTeX para facilitar digitação
%% 
%% =======================================================================================

\documentclass[pdftex,12pt,oneside,a4paper,english,french,spanish,brazil]{abntex2}


% ----------------------------------------------------------------------------------------
% - PACOTES
  \usepackage{sty/configs}                  % Pacotes utilizados
  \usepackage{sty/metadados_rt}             % Dados do órgão / cliente
  \usepackage{sty/dados_ccsec}               % Dados do CCSEC
  \usepackage{sty/vocabulario}              % Vocabulário específico do projeto
  \usepackage{sty/leiaute_rt}               % Definições de layout do documento
  \usepackage{ragged2e}
  \usepackage{helvet}
  \usepackage{titlesec}

\renewcommand{\chapnumfont}{\sffamily\normalsize\bfseries\sffamily}
\renewcommand{\chaptitlefont}{\sffamily\normalsize\bfseries\sffamily}

\titleformat{\section}{\sffamily\normalsize\bfseries}{}{0pt}{\thesection\space}
\titleformat{\subsection}{\sffamily\normalsize}{}{0pt}{\thesubsection\space}
\titleformat{\subsubsection}{\sffamily\normalsize\itshape}{}{0pt}{\thesubsubsection\space}


% ----------------------------------------------------------------------------------------

% ----------------------------------------------------------------------------------------
% - INÍCIO DO DOCUMENTO
  \begin{document}
      \begin{sloppypar}
          
          % Seleciona o idioma do documento (conforme pacotes do babel)
            %\selectlanguage{english}
            \selectlanguage{brazil}
            
          % Retira espaço extra obsoleto entre as frases.
            \frenchspacing
          
          % ------------------------------------------------------------------------------
          % - ESTRUTURA DO DOCUMENTO
          % ------------------------------------------------------------------------------
          % ------------------------------------------------------------------------------
          % - PRÉ-TEXTUAL
          % ------------------------------------------------------------------------------
          \pretextual

          % Capa
            \input{pre-textual/capa}
            \cleardoublepage
          
          % Revisões
            \input{pre-textual/revisoes}
            \cleardoublepage

          % Filiação CPAI
            \input{pre-textual/filiacao_ccsec}
            \cleardoublepage

          % Inserir a ficha bibliografica
            \input{pre-textual/ficha_catalografica}

            \begin{resumo}[RESUMO] 
                Resumo em português 
                
                \vspace{\onelineskip} 
                \noindent 
                
                
                \textbf{Palavras-chaves}:latex.abntex.editoraçãodetexto. 
            \end{resumo}

            \begin{resumo}[ABSTRACT] 
                Resumo em português 
                
                \vspace{\onelineskip} 
                \noindent 
                
                
                \textbf{Palavras-chaves}:latex, abntex, editoraçãodetexto. 
            \end{resumo}

          % Sumarios
            \renewcommand\contentsname{SUMÁRIO}

            \addtocontents{toc}{\protect\setlength{\protect\parskip}{10pt}} % Ajuste o '10pt' conforme necessário
            \hypersetup{
                colorlinks=true,       % Se true, os links são coloridos. Se false, os links são colocados em caixas.
                linkcolor=black,       % Cor dos links internos (ex. figuras, tabelas, seções)
                citecolor=black,       % Cor dos links para bibliografia
                filecolor=black,       % Cor dos links para arquivos
                urlcolor=black,        % Cor dos links para URLs externas
                bookmarksopen=true,    % Mostra os marcadores (ou bookmarks) do PDF quando aberto
            }
            
            \tableofcontents*
            
            \cleardoublepage
            
            \pagenumbering{roman}
            \renewcommand\listfigurename{LISTA DE FIGURAS}
            \listoffigures 
            \cleardoublepage

            \renewcommand\listtablename{LISTA DE TABELAS}
            \listoftables

            \cleardoublepage

                    
          % Lista de abreviaturas e siglas
            
            \include{pre-textual/siglas}

            \pagenumbering{arabic}
            
          % ------------------------------------------------------------------------------
          % - TEXTUAL
          % ------------------------------------------------------------------------------
          \textual
          
          % Introdução
            \chapter{INTRODUÇÃO}\label{cap:introducao}
            \input{textual/rt_introducao}

          % Capítulo
            \chapter{REFERÊNCIAS DE COMANDOS}\label{cap:refcomandos}
            \input{textual/refcomandos}

          % Capítulo
            \chapter{TÍTULO DESTE CAPÍTULO}\label{cap:capitulo1}
            \input{textual/rt_cap1}
          
          % Capítulo
            \chapter{ESTE CAPÍTULO TAMBÉM PRECISA DE TÍTULO}\label{cap:capitulo2}
            \input{textual/rt_cap2}
          
          % Considerações Finais
            \chapter{CONSIDERAÇÕES FINAIS}\label{cap:consideracoesfinais}
            \input{textual/rt_conclusao}
          
          % Fechamento
            \input{textual/fechamento}

          % ------------------------------------------------------------------------------
          % - PÓS-TEXTUAL
          % ------------------------------------------------------------------------------
          \postextual          
        
          % Referências bibliográficas
            \cleardoublepage
            \renewcommand{\bibname}{REFERÊNCIAS}
            \bibliography{bib/rt}
        
          % Glossário
            %\glossary              % Consulte o manual da classe abntex2 para orientações sobre o glossário.

          % Apêndices
            %\part*{Apêndices}
            %\begin{apendicesenv}

                %\chapter{Título do Apêndice}\label{apend:apendice1}
                %\input{pos-textual/apend01}

            %\end{apendicesenv}
                
          % Anexos
            \part*{Anexos}
            \begin{anexosenv}

                \chapter{Qual o título desse anexo?}\label{anex:anex01}
                \input{pos-textual/anex01}

            \end{anexosenv}
          
          % Índice remissimo
            % \printindex

    \end{sloppypar}
\end{document}
